% Options for packages loaded elsewhere
\PassOptionsToPackage{unicode}{hyperref}
\PassOptionsToPackage{hyphens}{url}
\documentclass[
  a4paper,
]{article}
\usepackage{xcolor}
\usepackage[margin=25mm]{geometry}
\usepackage{amsmath,amssymb}
\setcounter{secnumdepth}{-\maxdimen} % remove section numbering
\usepackage{iftex}
\ifPDFTeX
  \usepackage[T1]{fontenc}
  \usepackage[utf8]{inputenc}
  \usepackage{textcomp} % provide euro and other symbols
\else % if luatex or xetex
  \usepackage{unicode-math} % this also loads fontspec
  \defaultfontfeatures{Scale=MatchLowercase}
  \defaultfontfeatures[\rmfamily]{Ligatures=TeX,Scale=1}
\fi
\usepackage{lmodern}
\ifPDFTeX\else
  % xetex/luatex font selection
\fi
% Use upquote if available, for straight quotes in verbatim environments
\IfFileExists{upquote.sty}{\usepackage{upquote}}{}
\IfFileExists{microtype.sty}{% use microtype if available
  \usepackage[]{microtype}
  \UseMicrotypeSet[protrusion]{basicmath} % disable protrusion for tt fonts
}{}
\makeatletter
\@ifundefined{KOMAClassName}{% if non-KOMA class
  \IfFileExists{parskip.sty}{%
    \usepackage{parskip}
  }{% else
    \setlength{\parindent}{0pt}
    \setlength{\parskip}{6pt plus 2pt minus 1pt}}
}{% if KOMA class
  \KOMAoptions{parskip=half}}
\makeatother
\usepackage{longtable,booktabs,array}
\newcounter{none} % for unnumbered tables
\usepackage{calc} % for calculating minipage widths
% Correct order of tables after \paragraph or \subparagraph
\usepackage{etoolbox}
\makeatletter
\patchcmd\longtable{\par}{\if@noskipsec\mbox{}\fi\par}{}{}
\makeatother
% Allow footnotes in longtable head/foot
\IfFileExists{footnotehyper.sty}{\usepackage{footnotehyper}}{\usepackage{footnote}}
\makesavenoteenv{longtable}
\usepackage{graphicx}
\makeatletter
\newsavebox\pandoc@box
\newcommand*\pandocbounded[1]{% scales image to fit in text height/width
  \sbox\pandoc@box{#1}%
  \Gscale@div\@tempa{\textheight}{\dimexpr\ht\pandoc@box+\dp\pandoc@box\relax}%
  \Gscale@div\@tempb{\linewidth}{\wd\pandoc@box}%
  \ifdim\@tempb\p@<\@tempa\p@\let\@tempa\@tempb\fi% select the smaller of both
  \ifdim\@tempa\p@<\p@\scalebox{\@tempa}{\usebox\pandoc@box}%
  \else\usebox{\pandoc@box}%
  \fi%
}
% Set default figure placement to htbp
\def\fps@figure{htbp}
\makeatother
\setlength{\emergencystretch}{3em} % prevent overfull lines
\providecommand{\tightlist}{%
  \setlength{\itemsep}{0pt}\setlength{\parskip}{0pt}}
\usepackage{bookmark}
\IfFileExists{xurl.sty}{\usepackage{xurl}}{} % add URL line breaks if available
\urlstyle{same}
\hypersetup{
  pdftitle={The Third Thought Experiment: Reading Out the State of Two Bodies a{,} b from the Orbit of Two Celestial Bodies a{,} b - Obtaining the Mass Ratio{,} the Relative Distance and the Relative Time from Orbit and Gravitational-Wave Readouts{,} with Two Values of Undetermined Meaning and the Knowledge G = c = M = 1 Alone},
  pdfauthor={Noriaki Kihara (ORCID 0009-0004-6753-4020)},
  hidelinks,
  pdfcreator={LaTeX via pandoc}}

\title{The Third Thought Experiment: Reading Out the State of Two Bodies
a, b from the Orbit of Two Celestial Bodies a, b - Obtaining the Mass
Ratio, the Relative Distance and the Relative Time from Orbit and
Gravitational-Wave Readouts, with Two Values of Undetermined Meaning and
the Knowledge G = c = M = 1 Alone}
\author{Noriaki Kihara (ORCID 0009-0004-6753-4020)}
\date{2026-09-23 / v1.5 / Version DOI 10.5281/zenodo.22909661}

\sloppy
\emergencystretch=3em
\begin{document}
\maketitle

\textbf{Author:} Noriaki Kihara\\
\textbf{ORCID:} 0009-0004-6753-4020\\
\textbf{Version DOI:} 10.5281/zenodo.22909661\\
\textbf{Concept DOI:} 10.5281/zenodo.22909660\\
\textbf{Version:} v1.5\\
\textbf{Date:} 2026-09-23

\begin{center}\rule{0.5\linewidth}{0.5pt}\end{center}

\subsection{Abstract}\label{abstract}

The first thought experiment {[}1{]} and the second thought experiment
{[}2{]} started from two nameless values \(a, b\) and derived
Kepler-type orbits and an inverse-square force. On the other hand,
{[}2{]} concluded that the only quantity readable from the relative
motion of two bodies is the single coupling \(\mu_{12}\), and that the
mass ratio cannot be written with two values. This paper starts from a
doubt about that conclusion and re-examines it along a route different
from the previous ones. The orbits of two celestial bodies were computed
with the exact Newtonian equations and with general relativity (the
post-Newtonian approximation), and for the latter a numerical experiment
tested whether the mass ratio, the relative distance and the relative
time can be read out from the observed orbit.

The only condition imposed on the state \(\Psi\) is that distinguishing
a state from a relation requires at least two values of undetermined
meaning (P1). What \(a, b\) are (the semi-major and semi-minor axes of
the orbit, or the masses of the two bodies) is zero knowledge for the
observer. The readout side holds the knowledge \(G=c=M=1\) (\(M\) the
total mass) and reads the observational data as the motion of two
bodies. ``Two bodies'' is not a condition imposed on \(\Psi\); it is the
way the readout side reads the data.

In numerical experiment 01, shape point sequences of the relative orbit
- circle, ellipse, parabola and hyperbola - were generated as exact
solutions of the Newtonian equations, with the information of time,
position and orientation removed. In numerical experiment 02, orbits in
a strong field with \(a\) from 30 to 120 (\(b\) from 27 to 112) were
computed in the post-Newtonian approximation including radiation
reaction (1PN, 2PN, 2.5PN). Two kinds of data obtained by a distant,
light observer \(C\) were produced: the positions of the two celestial
bodies on the celestial sphere seen along the normal of the orbital
plane (observation 1), and the two polarizations of the gravitational
wave (observation 2). The distance, the conversion of \(M\) into
seconds, the origin of time, the orientation of the celestial sphere,
and which of the two bodies is the heavier were all hidden. The readout
used the observational data only, and the truth was read only in the
final cross-check.

For all three data sets, with mass ratios 1, 2 and 4, the mass ratio,
the orbital shape \(a, b\), the conversion \(T_M\) of \(M\) into
seconds, the distance \(D\), the relative distance \(r(t)\), the elapsed
time and the difference of the proper times of the two bodies were read
out. Most relative errors are \(10^{-4}\) to \(10^{-3}\), and apart from
the mass ratio of D2 read from the gravitational wave alone by way of
\(\nu\) (\(6\times10^{-2}\)) the largest is \(4.7\times10^{-3}\). The
noise added to the positions of observation 1 and to the amplitude of
observation 2 has a standard deviation of \(1/1000\) of the reference
quantity in each case. From observation 1 it was also possible to decide
correctly which of the two bodies is the heavier. In a strong field, if
the orbit in the two-dimensional plane can be read out, then the mass
ratio, the relative distance and the relative time can be read out. If
the observer \(C\) itself is counted as a state, the system is extended
to three states; this paper treats \(C\) as a test observer, and the
dynamics of the three relations is the subject of the next thought
experiment.

\textbf{Keywords:} two-body problem, readout, mass ratio, post-Newtonian
approximation, gravitational waves, proper time, zero knowledge, thought
experiment

\begin{center}\rule{0.5\linewidth}{0.5pt}\end{center}

\subsection{0. Position of this paper}\label{position-of-this-paper}

This paper is the third thought experiment under the design document
{[}0{]}. {[}1{]} and {[}2{]} started from two nameless values and a law,
and derived orbits from the side of the law. This paper runs in the
opposite direction: it looks at what can be read out about the system
from data obtained by observing an orbit.

The labels are the same as in {[}1, 2{]}. Only those used in this paper
are listed.

{\def\LTcaptype{none} % do not increment counter
\begin{longtable}[]{@{}
  >{\raggedright\arraybackslash}p{(\linewidth - 2\tabcolsep) * \real{0.5000}}
  >{\raggedright\arraybackslash}p{(\linewidth - 2\tabcolsep) * \real{0.5000}}@{}}
\toprule\noalign{}
\begin{minipage}[b]{\linewidth}\raggedright
Label
\end{minipage} & \begin{minipage}[b]{\linewidth}\raggedright
Meaning
\end{minipage} \\
\midrule\noalign{}
\endhead
\bottomrule\noalign{}
\endlastfoot
{[}Principle{]} & Philosophical principle \\
{[}Design assumption{]} & Assumption of the preliminary design
(removable, intended to be removed) \\
{[}Numerical{]} & Confirmed by numerical experiment \\
{[}Cross-check{]} & Correspondence with known physics; not a
derivation \\
\end{longtable}
}

\textbf{List of inputs (nothing else is used)}

{\def\LTcaptype{none} % do not increment counter
\begin{longtable}[]{@{}
  >{\raggedright\arraybackslash}p{(\linewidth - 4\tabcolsep) * \real{0.3333}}
  >{\raggedright\arraybackslash}p{(\linewidth - 4\tabcolsep) * \real{0.3333}}
  >{\raggedright\arraybackslash}p{(\linewidth - 4\tabcolsep) * \real{0.3333}}@{}}
\toprule\noalign{}
\begin{minipage}[b]{\linewidth}\raggedright
Symbol
\end{minipage} & \begin{minipage}[b]{\linewidth}\raggedright
Label
\end{minipage} & \begin{minipage}[b]{\linewidth}\raggedright
Content
\end{minipage} \\
\midrule\noalign{}
\endhead
\bottomrule\noalign{}
\endlastfoot
P1 & {[}Principle{]} & Distinguishing a state from a relation requires
at least two values of undetermined meaning (Section 2.1 of {[}1{]}). No
internal ground is posited for deciding which is the state and which is
the relation \\
P2 & {[}Principle{]} & A state alone is unreadable. What is readable is
only what passes through a readout (an interaction) (Section 5 of
{[}0{]}) \\
D7 & {[}Design assumption{]} & The readout side interprets the
observational data in the units \(G=c=M=1\) (\(M\) the total mass) \\
D8 & {[}Design assumption{]} & The readout side reads the observational
data as the motion of two bodies, using a known model of gravity and a
known map of the observation: the equations of relative motion with the
2PN conservative terms and the 2.5PN radiation reaction, the 1PN
center-of-mass relation, the lowest-order (quadrupole) gravitational
waveform, and the 1PN rate of proper time. This model has no independent
term of self-interaction (Section 2.2) \\
D9 & {[}Design assumption{]} & The observer \(C\) is light enough not to
affect the motion of the two bodies, is far away, and looks along the
normal of the orbital plane. It has coordinates on the celestial sphere
and a clock. It reads the two celestial bodies as point masses without
spin, and neglects the Doppler effect and the difference of the light
travel time. During the observation, the celestial coordinate system of
\(C\) is treated as a fixed inertial reference, and the proper motion of
the center of mass of the two bodies, the parallax, the acceleration of
\(C\) itself, and a time-dependent rotation of the coordinates are
neglected. The resolution of the observation is as given in Section
4.4 \\
\end{longtable}
}

What the readout side does not use: the values of \(a, b\) on the
generating side and the truth, the way the data were made (the random
seed and the assignment of the conditions), and the hidden conditions of
the observation (the distance, the conversion of \(M\) into seconds, the
origin of time, the rotation, reflection and origin of the celestial
sphere, and which of the two bodies is the heavier). The design
assumptions of {[}1, 2{]} (the invertible linear map D2, the product
operation D3, the three-dimensional readout space D5, the area clock D6)
are not used either. The only law used in this paper is D8, which the
readout side brings in as known.

The numbering of the design assumptions continues from D1 through D4 of
{[}1{]} and D5, D6 of {[}2{]}, and starts here at D7.

\begin{center}\rule{0.5\linewidth}{0.5pt}\end{center}

\subsection{1. The question}\label{the-question}

Section 8.2 of {[}2{]} (derivation T21b) states the following. With two
values, the mass ratio and the charge ratio cannot be written. The
reason is that the only quantity readable from the relative motion of
two bodies is the single coupling \(\mu_{12}\), from which the
individual ratios cannot be recovered.

A doubt about this conclusion remained. The question of this paper is
the following.

\begin{center}
\fbox{\parbox{0.92\linewidth}{\centering When a system written with two values of undetermined meaning is read with the knowledge $G=c=M=1$, can the mass ratio, the relative distance and the relative time be read out from the orbit?}}
\end{center}

\begin{center}\rule{0.5\linewidth}{0.5pt}\end{center}

\subsection{2. Setting}\label{setting}

\subsubsection{\texorpdfstring{2.1 The state \(\Psi\) and the two
values}{2.1 The state \textbackslash Psi and the two values}}\label{the-state-psi-and-the-two-values}

The only condition imposed on the state \(\Psi\) is P1. P1 is carried
over from {[}1{]}: distinguishing a state from a relation requires at
least two values of undetermined meaning. In the generation of this
paper, the two values \(a, b\) were taken to be the semi-major and
semi-minor axes of the relative orbit at the start (in units of \(M\)).
This is one way of writing on the generating side, and it does not mean
that P1 limits the degrees of freedom of \(\Psi\) to two.

What \(a, b\) are is zero knowledge for the observer. The readout does
not use the values of \(a, b\) on the generating side or the truth. What
the readout side holds is the known two-body model (D8) and the
coordinates in which that model is written. The definition of those
coordinates is the same as in the generation, but this is not knowledge
of the input of the generation (Section 7.1). The observational data are
determined by the configuration of the two bodies (positions and
velocities) alone, and do not depend on what names were given to the
input. If the same configuration is produced by another way of writing,
the observational data are the same and the values read out are the
same. For instance, one may take \(a, b\) to be the masses of the two
bodies and give the size and the shape of the orbit separately. Even in
the generation of this paper, the mass ratio is not put into \(a, b\);
it is given as a separate input.

\subsubsection{2.2 The knowledge on the readout
side}\label{the-knowledge-on-the-readout-side}

\textbf{The interpretation \(G=c=M=1\) (D7).} {[}Cross-check{]} In
Newtonian mechanics, multiplying the size of the orbit by \(\lambda\)
and the time by \(\lambda^{3/2}\) leaves the equations of motion
unchanged (mechanical similarity {[}4{]}). For this reason, \(G=M=1\)
alone provides no scale by which to compare the sizes of orbits. What
can be read from the shape is only the eccentricity \(e\) and the time
up to its unit.

Adding \(c=1\) makes the gravitational length \(GM/c^2\) a unit. The
dimensionless quantities that fix the orbit of two bodies become three:
the eccentricity \(e\), the field strength \(M/p\) (\(p\) the semi-latus
rectum), and the symmetric mass ratio \(\nu=m_Am_B/M^2\). To hold the
knowledge \(G=c=M=1\) is to be able to read the observed size and period
of the orbit in these units. When \(a, b\) are in the range of a few
tens, the field is strong enough for the periastron advance and the
effect of gravitational radiation to appear in the orbit.

\textbf{Reading as two bodies (D8).} The readout side reads the
observational data as two point masses moving under gravity. The law
used there is the equations of motion of the post-Newtonian
approximation including radiation reaction {[}9, 10{]}, the same as the
one used to generate the data. This law describes the relative motion of
two point masses. In the standard post-Newtonian point-mass model the
self-field singularity is regularized, and the two-body equations of
motion up to 2.5PN are finally written with the two masses {[}10{]}.
Accordingly, D8 does not add to the law of the readout an independent
term corresponding to the self-interaction of Section 6 of {[}0{]} (the
2.5PN radiation reaction is included as the reaction of the
gravitational wave emitted by the system). This is not a derivation from
Section 6 of {[}0{]}; it is a design assumption made explicit in
accordance with Section 12 of {[}0{]}. ``Two bodies'' is not a condition
imposed on \(\Psi\); it is this way of reading. The information obtained
by the readout is not equal to what \(\Psi\) holds, and the readout is
not a map of the whole of \(\Psi\).

\textbf{Observation along the normal (D9).} A tilted circle looks like
an ellipse. To read the shape of an orbit in a two-dimensional plane
presupposes that one is looking along the normal of that plane (or that
the tilt is known). This paper makes that presupposition explicit as D9.
D9 does not deny ``observation \(\equiv\) readout accompanied by
interaction'' of Section 5 of {[}0{]}; it is the approximation of
neglecting the dynamical back-reaction of the interaction accompanying
the observation on the motion of the two bodies.

If the observer \(C\) too is counted as one state, having the degrees of
freedom of position, orientation and clock, the system is extended to
three states: the two celestial bodies \(A, B\) and \(C\). What matters
in this paper, however, is not that the number of states becomes three,
but that there is a hierarchy of geometry among the three relations
\((A,B)\), \((A,C)\), \((B,C)\). \(C\) is placed far enough away, nearly
along the normal of the plane of the relative orbit of the two bodies,
and the viewing angle \(\theta_{AB}^{(C)}\) of the two bodies seen from
\(C\) is small enough (at most \(2\times10^{-4}\) rad in the setting of
this paper). On the other hand, the angle \(\theta_{BC}^{(A)}\) between
the directions of \(B\) and \(C\) seen from \(A\), and the angle
\(\theta_{AC}^{(B)}\) between the directions of \(A\) and \(C\) seen
from \(B\), are nearly \(\pi/2\). That is,

\[
\theta_{AB}^{(C)}\ll\theta_{AC}^{(B)},\ \theta_{BC}^{(A)}
\]

By this configuration, \(C\) can read out the relation of \(A, B\)
against the common celestial coordinates and clock. On the other hand,
this paper treats \(C\) as a test observer, light enough and far enough
away, and neglects the back-reaction of the observation by \(C\) on the
motion of the two celestial bodies. This is not a claim that the
back-reaction is exactly zero but a hypothesis, and its validity is left
as a task for later research. What this paper treats ends here. Treating
\(A, B, C\) all as states internal to the system, and asking what
self-consistent dynamics the three relations make, is the subject of the
next thought experiment, and the entry point to the three-body problem
in this research.

\subsubsection{2.3 The quantities to be read
out}\label{the-quantities-to-be-read-out}

The quantities seen from the observer \(C\) are the following.

\begin{itemize}
\tightlist
\item
  The mass ratio \(q=m_A/m_B\ (\ge 1)\) and which of the two bodies is
  the heavier. The symmetric mass ratio \(\nu\)
\item
  The shape of the orbit \(a, b\) (in units of \(M\))
\item
  The conversion \(T_M=GM/c^3\) {[}s{]} of \(M\) into seconds, and the
  distance \(D\) (in units of \(M\))
\item
  The relative distance \(r(t)\) and the elapsed time \(t\) (both in
  units of \(M\)). Here \(r\) is the coordinate separation
  \(|\mathbf x_A-\mathbf x_B|\) of the two bodies in the harmonic
  coordinates adopted in D8, and not the proper distance itself, which
  is coordinate independent
\end{itemize}

The quantities of each of the two bodies are the following (the heavier
is called \(A\)).

\begin{itemize}
\tightlist
\item
  The proper times \(\tau_A(t)\), \(\tau_B(t)\) (in units of \(M\)).
  From the 1PN metric containing the field of the other body only,
  \(d\tau_A/dt=1-m_B/r-v_A^2/2\) with \(v_A=m_Bv\) (\(v\) the relative
  velocity) {[}15, 16{]}. Likewise for \(\tau_B\).
\item
  The clock-rate-converted distances \(d_A=r\,d\tau_A/dt\),
  \(d_B=r\,d\tau_B/dt\). These are the coordinate separation \(r\)
  converted by the rate of the proper clock {[}15{]} of each body. This
  is not the operational radar distance defined from the proper time of
  a round trip of light; the propagation of light including the Shapiro
  delay {[}17{]} is not computed.
\end{itemize}

\begin{center}\rule{0.5\linewidth}{0.5pt}\end{center}

\subsection{3. Numerical experiment 01: orbits from the Newtonian
equations}\label{numerical-experiment-01-orbits-from-the-newtonian-equations}

\subsubsection{3.1 Purpose and design}\label{purpose-and-design}

The Newtonian two-body problem is exactly solvable. Passing to the
reduced mass, the relative motion becomes a one-body central-force
problem, and the orbit is a conic \(r(\theta)=p/(1+e\cos\theta)\) with
the origin at a focus {[}3, 4, 5{]}. Only the total mass \(M\) enters
the equation of the relative orbit; the mass ratio does not.
{[}Cross-check{]}

In numerical experiment 01, shape point sequences containing no
information of time were made from this exact solution. They are data
for a blind test, to examine whether the type of the orbit, the two
degrees of freedom of the shape, the position and the orientation of the
focus, and the position in time can be estimated from the shape and the
ordering alone.

\begin{itemize}
\tightlist
\item
  \textbf{Generating formula}: with a standard coordinate system whose
  origin is the focus and whose \(+x\) direction is that of the
  periastron, \(\mathbf q(\theta)=r(\theta)(\cos\theta,\sin\theta)\) was
  made and placed in the data coordinates by a random rotation
  \(R(\varphi)\) and a translation \(O\)
  (\(P=O+R(\varphi)\,\mathbf q(\theta)\)).
\item
  \textbf{The two degrees of freedom of the shape}: \((p, e)\). Here
  \(p\) was a uniform random number between 0.5 and 3.0. For ellipses,
  \(a=p/(1-e^2)\) and \(b=p/\sqrt{1-e^2}\) were also saved as derived
  values.
\item
  \textbf{Cases}: circle (\(e=0\)), ellipse (\(e=0.02\), 0.6, 0.97), a
  partial arc of a long thin ellipse (\(e=0.97\)), parabola (\(e=1\)),
  hyperbola (\(e=1.03\), 1.8), eight cases in all. Each case was made
  with two ways of spacing, giving 16 data sets in total (318 to 1274
  points each).
\item
  \textbf{Removal of time}: the points were placed at irregular
  intervals based on arc length, unrelated to time. If they were spaced
  by an angle proportional to time, the information of time would remain
  in the density of the points.
\item
  \textbf{Removal of preconceptions}: the position of the focus, the
  direction of the major axis and the direction of travel were given at
  random. The correspondence between the data set number and the case
  was also shuffled at random.
\item
  \textbf{The truth of time}: since \(GM\) is not fixed by the shape,
  time was saved as the dimensionless time
  \(\tau=\sqrt{GM/p^3}\,(t-t_p)=\int_0^\theta d\theta'/(1+e\cos\theta')^2\)
  with \(p\) as the unit of length.
\end{itemize}

\subsubsection{3.2 Self-verification
{[}Numerical{]}}\label{self-verification-numerical}

The generated data satisfied all the self-verifications. Satisfaction of
the equation of the conic, recovery of the ordering by a
nearest-neighbour sort, agreement of the analytic formula for time with
the numerical integration, and so forth: the largest residual was below
\(1\times10^{-14}\,p\). Point sequences regenerated from the saved
parameters alone also agreed to a relative error of \(3\times10^{-15}\).

\subsubsection{3.3 What this computation shows
{[}Cross-check{]}}\label{what-this-computation-shows-cross-check}

\begin{itemize}
\tightlist
\item
  The shape of the relative orbit is fixed by two values (\(p, e\), or
  \(a, b\)), and the mass ratio does not enter.
\item
  Time is fixed only up to the unit \(\sqrt{p^3/GM}\) and an origin,
  because \(GM\) is not fixed by the shape alone.
\item
  Reading the shape itself rests on the presupposition (D9) of looking
  along the normal of the orbital plane.
\end{itemize}

In numerical experiment 01, the generation and the self-verification
were carried out. A readout from the shape point sequences is not
performed in this paper. The readout experiment is carried out in
numerical experiment 02 below.

\begin{figure}
\centering
\pandocbounded{\includegraphics[keepaspectratio,alt={Figure 1a. The input of numerical experiment 01 (the point sequence as seen by the solving side)}]{figs/overview_input_en.png}}
\caption{Figure 1a. The input of numerical experiment 01 (the point
sequence as seen by the solving side)}
\end{figure}

\begin{figure}
\centering
\pandocbounded{\includegraphics[keepaspectratio,alt={Figure 1b. The truth of numerical experiment 01 (focus, periastron direction, direction of travel, time)}]{figs/overview_truth_en.png}}
\caption{Figure 1b. The truth of numerical experiment 01 (focus,
periastron direction, direction of travel, time)}
\end{figure}

\begin{center}\rule{0.5\linewidth}{0.5pt}\end{center}

\subsection{4. Numerical experiment 02: orbits from general relativity,
and the observational
data}\label{numerical-experiment-02-orbits-from-general-relativity-and-the-observational-data}

\subsubsection{4.1 The equations of
motion}\label{the-equations-of-motion}

The relative motion was computed with the equations of motion of the
post-Newtonian approximation in harmonic coordinates. To the Newtonian
term were added the 1PN and 2PN conservative terms and the 2.5PN
radiation-reaction term. Of the expression for the relative acceleration
given by Kidder {[}9{]}, the part not containing spin was used (review
{[}10{]}).

The integration was done with the fourth-order Runge-Kutta method, with
the step proportional to the local dynamical time,
\(h=(2\pi/4000)\,r^{3/2}\). The setting was to stop if the relative
distance fell below \(8M\), but no stop occurred in any of the data
sets.

The individual positions of the two bodies were obtained from the 1PN
center-of-mass relation {[}9, 10{]}.

\[
\mathbf x_A=\Bigl(m_B+\tfrac12\nu\,\delta m\,(v^2-1/r)\Bigr)\mathbf x,\qquad
\mathbf x_B=\Bigl(-m_A+\tfrac12\nu\,\delta m\,(v^2-1/r)\Bigr)\mathbf x
\]

Here \(\mathbf x=\mathbf x_A-\mathbf x_B\) and \(\delta m=m_A-m_B\). The
proper times \(\tau_A\), \(\tau_B\) were integrated with the formulas of
Section 2.3.

\subsubsection{\texorpdfstring{4.2 Definition of
\(a, b\)}{4.2 Definition of a, b}}\label{definition-of-a-b}

\(a, b\) are the semi-major and semi-minor axes of the relative orbit at
the start (\(t=0\)). In the 2PN motion without radiation reaction, the
tangential velocity at the periastron was fixed by bisection so that the
turning points (periastron and apastron) fall exactly at \(a(1-e)\) and
\(a(1+e)\) (with \(e=\sqrt{1-b^2/a^2}\)).

In a strong field this velocity differs from the Newtonian value by up
to about 10 per cent (0.2624 against 0.2913 at \(a=30\)). \(a, b\) are
not the Newtonian osculating ellipse; they are quantities defined by the
turning points of the 2PN motion.

\subsubsection{4.3 The three data sets}\label{the-three-data-sets}

{\def\LTcaptype{none} % do not increment counter
\begin{longtable}[]{@{}
  >{\raggedright\arraybackslash}p{(\linewidth - 14\tabcolsep) * \real{0.1250}}
  >{\raggedright\arraybackslash}p{(\linewidth - 14\tabcolsep) * \real{0.1250}}
  >{\raggedright\arraybackslash}p{(\linewidth - 14\tabcolsep) * \real{0.1250}}
  >{\raggedright\arraybackslash}p{(\linewidth - 14\tabcolsep) * \real{0.1250}}
  >{\raggedright\arraybackslash}p{(\linewidth - 14\tabcolsep) * \real{0.1250}}
  >{\raggedright\arraybackslash}p{(\linewidth - 14\tabcolsep) * \real{0.1250}}
  >{\raggedright\arraybackslash}p{(\linewidth - 14\tabcolsep) * \real{0.1250}}
  >{\raggedright\arraybackslash}p{(\linewidth - 14\tabcolsep) * \real{0.1250}}@{}}
\toprule\noalign{}
\begin{minipage}[b]{\linewidth}\raggedright
Data set
\end{minipage} & \begin{minipage}[b]{\linewidth}\raggedright
\(a\)
\end{minipage} & \begin{minipage}[b]{\linewidth}\raggedright
\(b\)
\end{minipage} & \begin{minipage}[b]{\linewidth}\raggedright
\(e\)
\end{minipage} & \begin{minipage}[b]{\linewidth}\raggedright
semi-latus rectum \(p\)
\end{minipage} & \begin{minipage}[b]{\linewidth}\raggedright
mass ratio \(q\)
\end{minipage} & \begin{minipage}[b]{\linewidth}\raggedright
\(\nu\)
\end{minipage} & \begin{minipage}[b]{\linewidth}\raggedright
computed length
\end{minipage} \\
\midrule\noalign{}
\endhead
\bottomrule\noalign{}
\endlastfoot
D1 & 120 & 112 & 0.3590 & 104.5 & 4 & 0.1600 & 30 times \(T_0\) \\
D2 & 30 & 27 & 0.4359 & 24.3 & 1 & 0.2500 & 20 times \(T_0\) \\
D3 & 60 & 50 & 0.5528 & 41.7 & 2 & 0.2222 & 30 times \(T_0\) \\
\end{longtable}
}

Here \(T_0=2\pi a^{3/2}\) is the Newtonian period. The field is
strongest in D2, whose periastron distance is \(16.9M\). The
correspondence between the data set numbers and the conditions was
shuffled at random on the generating side.

\subsubsection{\texorpdfstring{4.4 The observer \(C\) and the
observational
data}{4.4 The observer C and the observational data}}\label{the-observer-c-and-the-observational-data}

The observer \(C\) is far away along the normal of the orbital plane and
does not affect the motion of the two bodies. The Doppler effect and the
difference of the light travel time from the two bodies were neglected.

\begin{itemize}
\tightlist
\item
  \textbf{Observation 1 (orbit observation)}: the positions of each of
  the two bodies on the celestial sphere (in nrad) and the time of the
  observer (in seconds). The sampling interval is \(T_0/48\), and the
  error of the position is \(1/1000\) of the size \(a\) of the orbit
  converted into an angle on the celestial sphere
  (\(\sigma=10^{-3}\,a/D\) rad).
\item
  \textbf{Observation 2 (gravitational-wave observation)}: the two
  polarizations \(h_+\), \(h_\times\) of the gravitational wave and the
  time of the observer. The sampling interval is \(T_0/192\). The
  waveform is the quadrupole formula
  \(h_{ij}=(4\nu/D)(v_iv_j-n_in_j/r)\) {[}11, 12, 14{]}, and the error
  is \(1/1000\) of the maximum amplitude.
\end{itemize}

The hidden conditions of the observation (not given to the readout side)
are the following.

\begin{itemize}
\tightlist
\item
  The distance \(D\): a log-uniform random number between
  \(3\times10^5\) and \(3\times10^6\,M\)
\item
  The conversion \(T_M\) of \(M\) into seconds: fixed by taking the
  total mass to be a uniform random number between 20 and 80 solar
  masses
\item
  The origin of time: a shift between 0 and 1000 s
\item
  The rotation angle of the celestial sphere, the presence or absence of
  a reflection, and a shift of the origin on the celestial sphere (up to
  \(\pm2\times10^5\) nrad)
\item
  Whether ``body 1'' and ``body 2'' of observation 1 are the heavier or
  the lighter
\end{itemize}

\begin{figure}
\centering
\pandocbounded{\includegraphics[keepaspectratio,alt={Figure 2. The observational data of numerical experiment 02 (D1)}]{figs/obs_D1_en.png}}
\caption{Figure 2. The observational data of numerical experiment 02
(D1)}
\end{figure}

D2 and D3 are shown in obs\_D2\_en.png and obs\_D3\_en.png in the same
folder.

\subsubsection{4.5 Self-verification
{[}Numerical{]}}\label{self-verification-numerical-1}

{\def\LTcaptype{none} % do not increment counter
\begin{longtable}[]{@{}
  >{\raggedright\arraybackslash}p{(\linewidth - 6\tabcolsep) * \real{0.2500}}
  >{\raggedright\arraybackslash}p{(\linewidth - 6\tabcolsep) * \real{0.2500}}
  >{\raggedright\arraybackslash}p{(\linewidth - 6\tabcolsep) * \real{0.2500}}
  >{\raggedright\arraybackslash}p{(\linewidth - 6\tabcolsep) * \real{0.2500}}@{}}
\toprule\noalign{}
\begin{minipage}[b]{\linewidth}\raggedright
Item
\end{minipage} & \begin{minipage}[b]{\linewidth}\raggedright
D1
\end{minipage} & \begin{minipage}[b]{\linewidth}\raggedright
D2
\end{minipage} & \begin{minipage}[b]{\linewidth}\raggedright
D3
\end{minipage} \\
\midrule\noalign{}
\endhead
\bottomrule\noalign{}
\endlastfoot
Periastron advance per revolution \(\div\) the lowest-order value
\(6\pi/p\) {[}6{]} & 1.032 & 1.218 & 1.092 \\
Rate of energy loss \(\div\) the lowest-order formula of Peters {[}12,
13{]} & 0.974 & 0.982 & 0.968 \\
Difference of position on refining the step \(\div\ a\) &
\(5.8\times10^{-12}\) & \(3.8\times10^{-12}\) & \(9.1\times10^{-12}\) \\
\end{longtable}
}

The departure of the periastron advance from the lowest order is larger
the stronger the field. This agrees with the size of the 2PN correction.
In every data set the 2PN energy was better conserved than the 1PN
energy.

The consistency of the post-Newtonian approximation was also checked.
When the field strength was varied (\(\nu=0.25\), \(e=0.436\), \(a=60\)
to 480), the violation of the conservation of the 1PN and 2PN energies
decreased as the 1.83rd and 3.95th powers of the field strength at the
periastron respectively (the pass criteria were above 1.5 and above 2.7;
asymptotically the second and third powers are expected).

\begin{figure}
\centering
\pandocbounded{\includegraphics[keepaspectratio,alt={Figure 3. The truth of numerical experiment 02 (D1: relative orbit, relative distance, difference of proper times)}]{figs/truth_D1_en.png}}
\caption{Figure 3. The truth of numerical experiment 02 (D1: relative
orbit, relative distance, difference of proper times)}
\end{figure}

D2 and D3 are shown in truth\_D2\_en.png and truth\_D3\_en.png in the
same folder.

\begin{center}\rule{0.5\linewidth}{0.5pt}\end{center}

\subsection{5. The readout method}\label{the-readout-method}

\subsubsection{5.1 Principles of the
procedure}\label{principles-of-the-procedure}

The readout uses the observational data (observation 1 and observation
2) only. The truth is read only in the cross-check of Section 6. What
the readout side holds is the knowledge D7 through D9.

\subsubsection{5.2 Reading the mass ratio from the center of mass
(observation
1)}\label{reading-the-mass-ratio-from-the-center-of-mass-observation-1}

Writing the positions of the two bodies on the celestial sphere as
\(\mathbf X_1\), \(\mathbf X_2\) and
\(\boldsymbol\rho=\mathbf X_1-\mathbf X_2\), the 1PN center-of-mass
relation gives the following.

\[
\mathbf X_1=\mathbf X_0+\bigl(\mu_2+\nu(\mu_1-\mu_2)f(t)\bigr)\boldsymbol\rho,\qquad f=\tfrac12\left(v^2-1/r\right)
\]

Here \(\mu_i\) is the mass fraction of body \(i\) and \(\mathbf X_0\) is
the position of the center of mass on the celestial sphere. First, with
\(f=0\) (the Newtonian approximation), \(\mathbf X_0\) and \(\mu_2\)
were obtained by least squares. Then \(f(t)\) was computed from the
orbit reconstructed in Section 5.4 and they were obtained again. If
\(\mu_2<1/2\), body 1 is the heavier.

\subsubsection{5.3 Features at each periastron
passage}\label{features-at-each-periastron-passage}

\begin{itemize}
\tightlist
\item
  \textbf{Observation 1}: the distance and the angle of the relative
  position \(\boldsymbol\rho\) are used. At each periastron, a conic
  with periastron advance, \(1/r=A+B\cos\kappa\phi+C\sin\kappa\phi\),
  was fitted locally, and the time of the periastron, the direction of
  the periastron, the periastron distance and the following apastron
  distance were obtained. Here \(\kappa=1/(1+k)\), with \(k\) the
  fractional periastron advance per radial period.
\item
  \textbf{Observation 2}: the complex amplitude \(H=h_+-ih_\times\) is
  used. The time of a maximum of \(|H|\), the phase at that time, the
  maximum value and the following minimum value were obtained. For the
  quadrupole waveform, \(|H|\) is maximal at the periastron and minimal
  at the apastron, and the phase advances by \(4\pi(1+k)\) per
  periastron. The oscillation of the gravitational field itself plays
  the role of a clock.
\end{itemize}

\subsubsection{5.4 Fitting the sequence of model
features}\label{fitting-the-sequence-of-model-features}

A model orbit is computed with the equations of motion of D8, sampled at
the same times as the observation, and features are extracted by the
same procedure as in Section 5.3. These are fitted to the sequence of
features of the observation. In this way the bias originating in the way
the features are extracted cancels between the observation and the
model. The step of the model computation was 500 per revolution.

\begin{itemize}
\tightlist
\item
  \textbf{Nonlinear parameters}: \(a\), \(e\). For observation 1,
  \(\nu\) was fixed to the value from the center of mass of Section 5.2;
  for observation 2 it was free. A separate fit in which \(\nu\) was
  also free was performed for observation 1, and \(\nu\) was read from
  the decrease of the period.
\item
  \textbf{Linear parameters}: the origin of time, \(T_M\), the origin of
  the direction (phase), and the scale (nrad/\(M\) for observation 1,
  \(1/D\) for observation 2). These were eliminated by least squares
  {[}24{]}.
\item
  \textbf{Fitting}: the Levenberg-Marquardt method {[}22, 23{]} was
  used. The error of each feature was estimated from the residuals, and
  1\(\sigma\) was obtained from the covariance over all parameters.
\end{itemize}

The main clue that fixes each quantity is the following.
{[}Cross-check{]}

{\def\LTcaptype{none} % do not increment counter
\begin{longtable}[]{@{}
  >{\raggedright\arraybackslash}p{(\linewidth - 2\tabcolsep) * \real{0.5000}}
  >{\raggedright\arraybackslash}p{(\linewidth - 2\tabcolsep) * \real{0.5000}}@{}}
\toprule\noalign{}
\begin{minipage}[b]{\linewidth}\raggedright
Quantity read out
\end{minipage} & \begin{minipage}[b]{\linewidth}\raggedright
Clue
\end{minipage} \\
\midrule\noalign{}
\endhead
\bottomrule\noalign{}
\endlastfoot
\(a\) (in units of \(M\)) & the periastron advance per radial period
{[}6, 7, 8{]} \\
\(e\) & the ratio of the periastron and apastron distances (for
observation 2, the ratio of the maximum and minimum of the amplitude) \\
\(\nu\) & the decrease of the radial period (gravitational radiation)
{[}12, 13{]} \\
\(T_M\) & the radial period in seconds \(\div\) the period of the model
(in units of \(M\)) \\
\(D\) & the size in angle \(\div\) the size in units of \(M\)
(observation 1), the amplitude (observation 2) {[}21{]} \\
the mass ratio and which is the heavier & the individual positions of
the two bodies and the center of mass (observation 1) \\
\end{longtable}
}

\subsubsection{5.5 Reconstruction}\label{reconstruction}

With the parameters read out, the model orbit was rebuilt, and the
relative distance \(r(t)\), the elapsed time \(t\), the proper times
\(\tau_A\), \(\tau_B\) and the clock-rate-converted distances \(d_A\),
\(d_B\) at the times of the observation were computed (the heavier is
called \(A\)). The relative distance was also computed by the direct
method of dividing the angle of observation 1 by the scale that was read
out.

\begin{center}\rule{0.5\linewidth}{0.5pt}\end{center}

\subsection{6. Results {[}Numerical{]}}\label{results-numerical}

\subsubsection{6.1 Errors of the quantities read
out}\label{errors-of-the-quantities-read-out}

The table gives the relative error
\(|\text{readout}-\text{truth}|/\text{truth}\).

{\def\LTcaptype{none} % do not increment counter
\begin{longtable}[]{@{}
  >{\raggedright\arraybackslash}p{(\linewidth - 8\tabcolsep) * \real{0.2000}}
  >{\raggedright\arraybackslash}p{(\linewidth - 8\tabcolsep) * \real{0.2000}}
  >{\raggedright\arraybackslash}p{(\linewidth - 8\tabcolsep) * \real{0.2000}}
  >{\raggedright\arraybackslash}p{(\linewidth - 8\tabcolsep) * \real{0.2000}}
  >{\raggedright\arraybackslash}p{(\linewidth - 8\tabcolsep) * \real{0.2000}}@{}}
\toprule\noalign{}
\begin{minipage}[b]{\linewidth}\raggedright
Quantity
\end{minipage} & \begin{minipage}[b]{\linewidth}\raggedright
Method
\end{minipage} & \begin{minipage}[b]{\linewidth}\raggedright
D1
\end{minipage} & \begin{minipage}[b]{\linewidth}\raggedright
D2
\end{minipage} & \begin{minipage}[b]{\linewidth}\raggedright
D3
\end{minipage} \\
\midrule\noalign{}
\endhead
\bottomrule\noalign{}
\endlastfoot
mass ratio \(q\) & center of mass (observation 1) & \(2.0\times10^{-4}\)
& \(1.9\times10^{-4}\) & \(3.2\times10^{-5}\) \\
mass ratio \(q\) & gravitational wave (from \(\nu\)) &
\(7.9\times10^{-4}\) & \(6.0\times10^{-2}\) & \(3.3\times10^{-3}\) \\
\(\nu\) & center of mass (observation 1) & \(1.2\times10^{-4}\) &
\(8.7\times10^{-9}\) & \(1.1\times10^{-5}\) \\
\(\nu\) & decrease of the period (observation 1) & \(4.3\times10^{-3}\)
& \(1.4\times10^{-4}\) & \(9.1\times10^{-4}\) \\
\(\nu\) & gravitational wave & \(4.8\times10^{-4}\) &
\(8.6\times10^{-4}\) & \(1.1\times10^{-3}\) \\
\(a\) & observation 1 / gravitational wave & \(4.4\times10^{-4}\) /
\(1.0\times10^{-3}\) & \(1.8\times10^{-4}\) / \(5.1\times10^{-5}\) &
\(2.4\times10^{-4}\) / \(6.3\times10^{-4}\) \\
\(b\) & observation 1 / gravitational wave & \(4.9\times10^{-4}\) /
\(8.9\times10^{-4}\) & \(1.5\times10^{-4}\) / \(7.6\times10^{-5}\) &
\(1.9\times10^{-4}\) / \(3.3\times10^{-4}\) \\
\(e\) & observation 1 / gravitational wave & \(3.2\times10^{-4}\) /
\(7.0\times10^{-4}\) & \(1.4\times10^{-4}\) / \(5.4\times10^{-4}\) &
\(1.3\times10^{-4}\) / \(6.8\times10^{-4}\) \\
\(T_M\) & observation 1 / gravitational wave & \(6.6\times10^{-4}\) /
\(1.5\times10^{-3}\) & \(3.1\times10^{-4}\) / \(4.5\times10^{-6}\) &
\(3.5\times10^{-4}\) / \(9.0\times10^{-4}\) \\
\(D\) & observation 1 / gravitational wave & \(3.5\times10^{-4}\) /
\(1.2\times10^{-3}\) & \(3.3\times10^{-4}\) / \(6.6\times10^{-4}\) &
\(3.7\times10^{-4}\) / \(7.9\times10^{-4}\) \\
\(\tau_A-\tau_B\) (final) & observation 1 / gravitational wave &
\(3.2\times10^{-4}\) / \(9.1\times10^{-4}\) & truth is 0 (*) &
\(5.9\times10^{-5}\) / \(4.7\times10^{-3}\) \\
\(d_A-d_B\) (mean) & observation 1 / gravitational wave &
\(1.2\times10^{-4}\) / \(4.0\times10^{-4}\) & truth is 0 (*) &
\(5.3\times10^{-5}\) / \(4.4\times10^{-3}\) \\
root mean square of \(r(t)\) & angle of observation 1 \(\div\) scale &
\(1.5\times10^{-3}\) & \(1.8\times10^{-3}\) & \(1.7\times10^{-3}\) \\
root mean square of \(r(t)\) & reconstruction (observation 1 /
gravitational wave) & \(4.6\times10^{-4}\) / \(1.1\times10^{-3}\) &
\(3.9\times10^{-4}\) / \(3.3\times10^{-4}\) & \(3.0\times10^{-4}\) /
\(7.8\times10^{-4}\) \\
\end{longtable}
}

(*) D2 has equal masses, so the truth of \(\tau_A-\tau_B\) and
\(d_A-d_B\) is 0. The readout from observation 1 gave
\(5.7\times10^{-6}\) of \(\tau_A\) and \(5.1\times10^{-6}\) of \(d_A\)
respectively. From the gravitational wave it gave \(1.8\times10^{-3}\)
of \(\tau_A\) and \(1.6\times10^{-3}\) of \(d_A\).

The error of the final value of the elapsed time (in units of \(M\)) was
\(6.6\times10^{-4}\) for D1, \(3.0\times10^{-4}\) for D2 and
\(3.6\times10^{-4}\) for D3 in observation 1. This agrees with the error
of \(T_M\).

\subsubsection{6.2 Examples of the values read
out}\label{examples-of-the-values-read-out}

{\def\LTcaptype{none} % do not increment counter
\begin{longtable}[]{@{}
  >{\raggedright\arraybackslash}p{(\linewidth - 6\tabcolsep) * \real{0.2500}}
  >{\raggedright\arraybackslash}p{(\linewidth - 6\tabcolsep) * \real{0.2500}}
  >{\raggedright\arraybackslash}p{(\linewidth - 6\tabcolsep) * \real{0.2500}}
  >{\raggedright\arraybackslash}p{(\linewidth - 6\tabcolsep) * \real{0.2500}}@{}}
\toprule\noalign{}
\begin{minipage}[b]{\linewidth}\raggedright
Data set
\end{minipage} & \begin{minipage}[b]{\linewidth}\raggedright
Quantity
\end{minipage} & \begin{minipage}[b]{\linewidth}\raggedright
Readout
\end{minipage} & \begin{minipage}[b]{\linewidth}\raggedright
Truth
\end{minipage} \\
\midrule\noalign{}
\endhead
\bottomrule\noalign{}
\endlastfoot
D1 & mass ratio \(q\) (center of mass) & 4.0008 & 4 \\
D1 & \(\tau_A-\tau_B\) (final, observation 1) & 1853.15 \(M\) & 1852.56
\(M\) \\
D1 & \(T_M\) (observation 1) & \(3.4320\times10^{-4}\) s &
\(3.4343\times10^{-4}\) s \\
D1 & \(D\) (observation 1) & \(8.740\times10^5\,M\) &
\(8.737\times10^5\,M\) \\
D2 & mass ratio \(q\) (center of mass) & 1.0002 & 1 \\
D2 & \(a\) / \(b\) (gravitational wave) & 30.002 / 26.998 & 30 / 27 \\
D3 & mass ratio \(q\) (center of mass) & 2.0001 & 2 \\
D3 & \(\tau_A-\tau_B\) (final, observation 1) & 778.85 \(M\) & 778.90
\(M\) \\
\end{longtable}
}

Which of the two bodies is the heavier was decided correctly for D1 and
D3 from the center of mass of observation 1. D2 has equal masses, so it
is not subject to the decision.

\subsubsection{6.3 Validity of the fit, and the remaining
bias}\label{validity-of-the-fit-and-the-remaining-bias}

In the fit, the scale of the error of each feature is fixed from the
root mean square of the same residuals (Section 5.4). For this reason
\(\chi^2\) per degree of freedom is nearly equal to \(N/(N-p)\), which
is fixed by the number of residuals \(N\) and the number of parameters
\(p\) alone (1.05 to 1.09), and is not an independent goodness-of-fit
test against the noise of the observation. The 1\(\sigma\) obtained from
the covariance and the pull are likewise local estimates based on these
residuals. The actual error against the estimated 1\(\sigma\) (the pull)
fell mostly within \(\pm2\), but since there is only one realization of
the noise, the coverage of the 1\(\sigma\) has not been checked.

The \(e\) read from the gravitational wave was displaced to the positive
side by \(5\times10^{-4}\) to \(7\times10^{-4}\) in all three data sets
(pull 2.5 to 3.2). This is the bias of obtaining the values of the
maxima and minima of the amplitude from noisy data, which is not
cancelled on the side of the noiseless model. Part of the displacement
of \(a\) and \(T_M\) read from the gravitational wave (pull about 2.8 in
D3) is linked to this.

The mass ratio obtained from the gravitational wave alone goes by way of
\(\nu\), so it cannot distinguish \(q\) from \(1/q\). Near equal masses
the variation of \(q\) with respect to \(\nu\) is large, so in D2 an
error of \(8.6\times10^{-4}\) in \(\nu\) gave \(q=1.06\). For the same
reason, in the reconstruction of D2 from the gravitational wave,
\(\nu=0.24979\) was treated as a mass difference of about 0.03, and
\(\tau_A-\tau_B\) became 35.9 \(M\) at the end (the truth is 0).

The fit of D2 with \(\nu\) free in observation 1 gave \(\nu=0.25003\),
slightly above the physical upper limit \(1/4\). In the readout, the
search range of \(\nu\) was taken to be 0.01 to 0.27 so that the search
does not stick at the boundary of equal masses, and the upper limit was
not imposed as a hard boundary. This value is treated as an estimate of
a coefficient near the boundary (Section 7.4).

\begin{figure}
\centering
\pandocbounded{\includegraphics[keepaspectratio,alt={Figure 4. The relative errors of the quantities read out, and the 1 sigma estimated by the readout}]{figs/readout_summary_en.png}}
\caption{Figure 4. The relative errors of the quantities read out, and
the 1 sigma estimated by the readout}
\end{figure}

\begin{figure}
\centering
\pandocbounded{\includegraphics[keepaspectratio,alt={Figure 5. Inside the readout: the advance of the direction of the periastron, the decrease of the radial period, the relative error of r(t), and tau\_A - tau\_B}]{figs/readout_details_en.png}}
\caption{Figure 5. Inside the readout: the advance of the direction of
the periastron, the decrease of the radial period, the relative error of
r(t), and tau\_A - tau\_B}
\end{figure}

\begin{center}\rule{0.5\linewidth}{0.5pt}\end{center}

\subsection{7. Discussion}\label{discussion}

\subsubsection{\texorpdfstring{7.1 The meaning of \(a, b\) is not needed
for the
readout}{7.1 The meaning of a, b is not needed for the readout}}\label{the-meaning-of-a-b-is-not-needed-for-the-readout}

The readout does not use the values of \(a, b\) on the generating side
or the files of the truth. What it used is the observational data and
the knowledge D7 through D9 only. The mass ratio was given without being
put into \(a, b\), and even so it was read out. Inside the readout, the
model was fitted in the coordinates \(a, e, \nu\). These coordinates use
the same definition as the generation, namely that the turning points of
the conservative 2PN motion are \(a(1\pm e)\), but they are coordinates
chosen by the readout side in order to write the model, and this is not
knowledge of the input of the generation. The optimum of the fit is
unchanged under a one-to-one change of coordinates, so fitting in other
coordinates should give the same mass ratio, relative distance and
relative time; a re-implementation in other coordinates is not carried
out in this paper. What this paper has confirmed is therefore that,
without knowing the values of \(a, b\) on the generating side, the
parameters of the model of D8 are fixed by the observational data.

With the reading as two bodies (D8), the sequence of features extracted
from the observation could be fitted consistently by the sequence of
features of a single model (the \(\chi^2\) per degree of freedom of
Section 6.3 is an internal diagnostic whose scale is fixed by the
residuals, not an independent goodness-of-fit test). This shows that
``the inverse problem of this paper closes under the reading as two
bodies'', and not that ``\(\Psi\) is two bodies''.

The information obtained by the readout is not equal to what \(\Psi\)
holds. The \(\nu\) readable from the relative motion alone is unchanged
under exchange of the two bodies. From the relative position of
observation 1 or from the gravitational wave alone, which is the heavier
cannot be read. It was possible to say which is the heavier only when
the positions of the two bodies were seen separately on the celestial
sphere of the observer \(C\). In D2, with equal masses, they cannot be
distinguished even then.

\subsubsection{7.2 The knowledge added to the
readout}\label{the-knowledge-added-to-the-readout}

The readout rested on the following knowledge.

\begin{itemize}
\tightlist
\item
  Interpreting in the units \(G=c=M=1\) (D7)
\item
  Reading as the motion of two bodies, using a known model and a known
  map of the observation (D8). The model is the one of the
  post-Newtonian approximation, the same as in the generation of the
  data
\item
  That the observer \(C\) looks along the normal, and is light and far
  away and has coordinates on the celestial sphere (D9). That the two
  bodies are read as point masses without spin, and that the Doppler
  effect and the difference of the light travel time are neglected
\end{itemize}

The meaning of \(a, b\) remained zero knowledge, and what the readout
needed was this knowledge alone.

\subsubsection{7.3 Relation to the conclusion of
{[}2{]}}\label{relation-to-the-conclusion-of-2}

Section 8.2 of {[}2{]} (derivation T21b) is a statement about the
relative motion of two bodies under an inverse-square coupling, and
within that range it still holds. The mass ratio could be read in this
paper because the material for the readout increased. One part is that
observation 1 contains the individual positions of the two bodies, which
is what Section 8.2 of {[}2{]} called ``additional relational
information allowing two accelerations to be compared against one and
the same inertial reference'', required for reading the mass ratio (the
structure treated in Section 8.3 of {[}2{]}, of reading the mass ratio
from the directed couplings of three interacting bodies, is not realized
in this paper; Section 2.2). The other part is that in a strong field
\(\nu\) enters the relative motion itself (it enters the equations of
motion from the 1PN term; it does not enter the lowest-order periastron
advance \(6\pi/p\), and appears in the 2PN term of the periastron
advance {[}8{]} and in the radiation reaction {[}12, 13{]}). This paper
does not contradict {[}2{]}. Also, in this paper the mass ratio is given
as an input of the generation, and whether the mass ratio can be written
with two values is not tested.

Reading masses from relativistic orbital effects (periastron advance,
delay of proper time, decrease of the period) is what has been done with
binary pulsars {[}18, 19{]} and with gravitational-wave observations
{[}20{]}, and the theoretical ground used in this paper is the same as
theirs. What is specific to this paper lies in the setting: the meaning
of \(a, b\) is zero knowledge, and the knowledge of the readout side is
limited to \(G=c=M=1\) and the reading as two bodies.

\subsubsection{7.4 Limits}\label{limits}

\begin{itemize}
\tightlist
\item
  The orbits of general relativity were computed in the post-Newtonian
  approximation (the 2PN conservative terms and the 2.5PN radiation
  reaction). This is not a full computation in general relativity. The
  field strength was limited to the range in which the periastron
  distance is \(16.9M\) or more.
\item
  The readout used the same equations of motion as the generation of the
  data. Moreover the readout program imports and uses, as they are, the
  functions for integration, interpolation and apastron search of the
  generating program. What this paper has confirmed is that, under the
  same known model and the same implementation, the parameters are fixed
  by the observed quantities, and that the readout procedure from the
  extraction of features to the inverse problem works. Robustness
  against model error of the equations of motion and against differences
  between independent implementations has not been tested.
\item
  The step of the model computation of the readout is 500 per
  revolution, coarser than the roughly 4000 of the generation. Redoing
  the readout of D2 with a step of 1000, the change of the parameters
  read out was at most \(8\times10^{-9}\) in relative terms, which is
  \(1\times10^{-5}\) times the estimated 1\(\sigma\) and sufficiently
  smaller than the error of the readout (stepsize\_check\_D2.csv in
  results).
\item
  The noise of each data set is one realization. The \(e\) read from the
  gravitational wave retains the bias of Section 6.3.
\item
  In the fits with \(\nu\) free, the search range was
  \(0.01\le\nu\le0.27\), and the physical upper limit \(\nu\le1/4\) was
  not imposed as a hard boundary. The free fit of D2 in observation 1
  gave \(\nu=0.25003\), slightly above the upper limit (Section 6.3).
\item
  The proper times were obtained in the 1PN approximation containing the
  field of the other body only. The clock-rate-converted distance \(d\)
  is the coordinate separation converted by the rate of the proper
  clock; it is not the operational radar distance and does not include
  the Shapiro delay. \(r\) is the coordinate separation in harmonic
  coordinates.
\item
  In numerical experiment 01, the generation and the self-verification
  were carried out. A readout from the shape point sequences was not
  carried out.
\item
  That the observation by the observer \(C\) does not affect the motion
  of the two celestial bodies is a hypothesis under the approximation of
  a test observer, with \(C\) placed far enough away along the normal of
  the plane of the relative orbit and the viewing angle small enough
  (D9, Section 2.2). Counting \(C\) as a state too, the system has three
  states, and the back-reaction of \(C\) and the dynamics of the three
  relations are left as tasks for the next thought experiment.
\end{itemize}

\begin{center}\rule{0.5\linewidth}{0.5pt}\end{center}

\subsection{8. Conclusion}\label{conclusion}

The only condition imposed on the state \(\Psi\) was that distinguishing
a state from a relation requires at least two values of undetermined
meaning (P1). What \(a, b\) are was taken to be zero knowledge, and the
readout side, holding the knowledge \(G=c=M=1\), read the observational
data as the motion of two bodies.

{[}Numerical{]} In this setting, from the orbit observation and the
gravitational-wave observation of two bodies in a strong field, the mass
ratio, the shape of the orbit, the conversion of \(M\) into seconds, the
distance, the relative distance, the elapsed time and the difference of
the proper times of the two bodies were read out. The noise added to the
observation is \(1/1000\) of the reference quantity, and most relative
errors are \(10^{-4}\) to \(10^{-3}\); apart from the mass ratio of D2
read from the gravitational wave alone by way of \(\nu\)
(\(6\times10^{-2}\), Section 6.3) the largest is \(4.7\times10^{-3}\).
Which of the two bodies is the heavier could also be decided from the
observation of the positions of the two bodies.

In this paper, ``the orbit in the two-dimensional plane can be read
out'' means that a sequence of positions of the two bodies on the
celestial sphere, time-stamped by the clock of the observer \(C\)
(observation 1), is obtained. A readout from the shape alone without
time (the data of numerical experiment 01) is not carried out in this
paper.

In a strong field, if the orbit in the two-dimensional plane can be read
out, then the mass ratio, the relative distance and the relative time
can be read out.

In this paper it was confirmed that information not obtainable from the
relative motion of the two celestial bodies alone can be read out by way
of the common celestial coordinates and clock of the observer \(C\). If
\(C\) itself is counted as a state having degrees of freedom, the system
is extended from two states to three, and the observational
configuration of this paper has the hierarchy of relations
\(\theta_{AB}^{(C)}\ll\theta_{AC}^{(B)},\theta_{BC}^{(A)}\) (Section
2.2). This paper treated \(C\) as a test observer, light enough and far
enough away, and neglected its back-reaction. Treating \(A, B, C\) all
as states internal to the system, and investigating the self-consistent
dynamics made by the three relations \((A,B)\), \((A,C)\), \((B,C)\), is
beyond the range of this paper. This is the next thought experiment, and
the entry point to the three-body problem in this research.

\begin{center}\rule{0.5\linewidth}{0.5pt}\end{center}

\subsection{Location of the data and
programs}\label{location-of-the-data-and-programs}

Everything is under the folder
anonymous\_vertex\_state\_generation\_geometry (Google Drive).

\begin{itemize}
\tightlist
\item
  \textbf{Numerical experiment 01}:
  conic\_orbit\_shape\_point\_generation\_v1/
  (generate\_conic\_orbit\_samples.py, reproduce\_from\_meta.py,
  plot\_conic\_orbit\_samples.py, README.md). The random seed is fixed
  (20260922), and the data and the figures can be regenerated by running
  the programs.
\item
  \textbf{Numerical experiment 02}:
  strong\_field\_binary\_orbit\_observation\_generation\_v1/
  (generate\_binary\_pn\_observations.py,
  plot\_binary\_pn\_observations.py, README.md, data/input, data/truth,
  figures). The random seed is 2026092202.
\item
  \textbf{Readout}:
  strong\_field\_binary\_orbit\_readout\_experiment\_v1/
  (readout\_binary\_pn.py, plot\_readout\_results.py, README.md,
  results/, figures/).
\end{itemize}

\begin{center}\rule{0.5\linewidth}{0.5pt}\end{center}

\subsection{Work record}\label{work-record}

\begin{itemize}
\tightlist
\item
  In the generation of numerical experiment 02, it was found that if the
  tangential velocity at the periastron is given by the Newtonian value,
  the turning points depart greatly from \(a(1\pm e)\) in a strong
  field. The definition was therefore changed to that of Section 4.2
  (the turning points of the 2PN motion).
\item
  In the first version of the readout, the estimated error was about one
  third of the actual error, too small, and there was also an error in
  the computation of the error of the distance \(D\). Moreover, a test
  on noiseless data showed that in D2, where the field is strongest, the
  extraction of the features itself has a bias (\(5\times10^{-4}\) of
  the radial period in the time of the maximum of the gravitational-wave
  amplitude). The method was changed to one in which the model orbit is
  also sampled at the same times as the observation and extracted by the
  same procedure, the ends of the fitting window were made to fall
  smoothly, and the computation of the covariance was corrected. Section
  6 gives the results after this improvement.
\item
  The bias of the \(e\) read from the gravitational wave (Section 6.3)
  remains after the improvement.
\item
  Points corrected in v1.1 (2026-09-23). ``In a strong field'' was added
  to the one-line conclusion (to make it compatible with Section 3.3).
  ``Two bodies'' was rewritten as ``not a condition imposed on
  \(\Psi\)'' rather than as a property of \(\Psi\). It was made explicit
  that the degrees of freedom of \(\Psi\) are not necessarily only two,
  and that the mass ratio was given as an input of the generation. The
  order at which \(\nu\) enters the equations of motion (from 1PN) was
  corrected. It was written that D8 has no term of self-interaction,
  that D9 makes explicit the celestial coordinates, the clock, the point
  masses and the neglect of the Doppler effect, and that D9 is an
  approximation from Section 5 of {[}0{]}. The names of the design
  assumptions of {[}1, 2{]} and the cited place in {[}2{]} (Section 8.2,
  T21b) were matched to the originals. The range of \(a, b\) in the
  abstract, the exception of the mass ratio from the gravitational wave
  alone, the unit of the error of observation 1, the wording of the pass
  criteria of the PN consistency check, the way the formula of Kidder is
  cited, and a missing label were corrected. The effect of the step of
  the model computation of the readout was checked with D2 (Section
  7.4).
\item
  Points corrected in v1.2 (2026-09-23, in reply to an external review).
  P1 was matched to the definition of {[}1{]} (two values of
  undetermined meaning). It was made explicit that what the readout does
  not use are the values of \(a, b\) on the generating side and the
  truth, and it was written that the coordinates of the model of the
  readout use the same definition as the generation and that the readout
  program shares functions with the generating program. It was made
  explicit that the \(\chi^2\) per degree of freedom is an internal
  diagnostic whose scale is fixed by the residuals (\(N/(N-p)\)) and not
  a goodness-of-fit test. ``Radar distance'' was changed to
  ``clock-rate-converted distance'', and it was written that \(r\) is
  the coordinate separation in harmonic coordinates. It was recorded
  that the free fit of \(\nu\) in D2 slightly exceeded the upper limit
  \(1/4\). The range of the errors in the abstract and the conclusion
  was matched to the actual distribution, and the meaning of ``the orbit
  can be read out'' was defined. The relation of D8 and D9 to Sections 5
  and 6 of {[}0{]} was rewritten as design assumptions.
\item
  Points added in v1.3 (2026-09-23). It was made explicit in Sections
  2.2 and 7.4 that including the observer \(C\) the system is strictly
  one of three states, that \(C\) is placed at nearly infinite distance
  along the normal and the hypothesis is made, under the approximation
  of a sufficiently narrow viewing angle, that the observation by \(C\)
  does not affect the states \(a, b\), and that the validity of this is
  a task for later research. It was also written that this point is
  consistent with Section 8.3 of {[}2{]} (what becomes readable with
  three bodies).
\item
  Points corrected in v1.4 (2026-09-23, in reply to a second review).
  That the observer \(C\) can be counted as a state having the degrees
  of freedom of position, orientation and clock, the hierarchy of
  viewing angles
  \(\theta_{AB}^{(C)}\ll\theta_{AC}^{(B)},\theta_{BC}^{(A)}\) among the
  three relations \((A,B)\), \((A,C)\), \((B,C)\), and that this paper
  treats \(C\) as a test observer while the dynamics of the three
  relations is the subject of the next thought experiment, were written
  in the abstract and in Sections 2.2, 7.4 and 8. The sentence of v1.3
  saying that this is consistent with Section 8.3 of {[}2{]} was
  reorganized as the additional relational information of Section 8.2 of
  {[}2{]} (T21b), and it was restated that the structure of three
  interacting bodies of Section 8.3 is not realized in this paper. The
  two celestial bodies are written \(A, B\) and the values \(a, b\),
  distinctly. What the readout does not use was changed to ``the values
  of \(a, b\) on the generating side and the truth'', and the use of the
  coordinates of the known model was made explicit in Sections 0 and 2.1
  as well. D8 lists the known model and the map of the observation
  (equations of motion, center-of-mass relation, quadrupole waveform,
  rate of proper time), and D9 adds that the celestial coordinate system
  is treated as a fixed inertial reference. P1 was unified with ``at
  least two''.
\item
  v1.5 (2026-09-23, final correction before publication). The symbol of
  the observer was changed from the lower-case \(c\) to the upper-case
  \(C\). The reason is that it collided with the speed of light \(c\) of
  \(G=c=M=1\), and did not correspond with the two celestial bodies
  \(A, B\). The symbol of the position of the center of mass on the
  celestial sphere in Section 5.2 was changed from \(\mathbf C\) to
  \(\mathbf X_0\) accordingly. The past lines of the work record were
  changed to the same symbols.
\end{itemize}

\begin{center}\rule{0.5\linewidth}{0.5pt}\end{center}

\subsection{References}\label{references}

\textbf{Self-citations}

{[}0{]} Noriaki Kihara, ``Designing a Research Project to Explore the
Emergence of Spacetime, Particles, and Interactions from States and
Relations: A Framework of Preliminary Design and Feasibility
Verification for Foundational-Physics Model Exploration under
Uncertainty'', v1.0 (2026). Concept DOI: 10.5281/zenodo.22851944
(Version DOI: 10.5281/zenodo.22851945)

{[}1{]} Noriaki Kihara, ``The First Thought Experiment: Tracing the
Equivalence Principle Down to Two Nameless Degrees of Freedom - Deriving
the Sign of the Sum of Squares, the Curvature Radius, and Three Systems
(Inertial, Uniformly Accelerated, Rotating) from the Conservation of an
Area Readout'', v1.0 (2026). Concept DOI: 10.5281/zenodo.22857952
(Version DOI: 10.5281/zenodo.22857953)

{[}2{]} Noriaki Kihara, ``The Second Thought Experiment: Deriving the
Coulomb-Type Inverse-Square Force from a Product Readout and an Area
Clock - Keeping Two Values and a Linear Law, Obtaining the Focus,
Attraction and Repulsion and Neutrality, and Fixing What Two Values
Cannot Write'', v1.2 (2026). Concept DOI: 10.5281/zenodo.22867336
(Version DOI: 10.5281/zenodo.22876596)

\textbf{The Newtonian two-body problem}

{[}3{]} I. Newton, \emph{Philosophiae Naturalis Principia Mathematica}
(London, 1687).

{[}4{]} L. D. Landau and E. M. Lifshitz, \emph{Mechanics}, 3rd ed.,
Course of Theoretical Physics Vol. 1 (Pergamon, Oxford, 1976).

{[}5{]} H. Goldstein, C. P. Poole, and J. L. Safko, \emph{Classical
Mechanics}, 3rd ed.~(Addison-Wesley, 2002).

\textbf{The two-body motion in general relativity}

{[}6{]} A. Einstein, ``Erklaerung der Perihelbewegung des Merkur aus der
allgemeinen Relativitaetstheorie,'' \emph{Sitzungsber. Preuss. Akad.
Wiss. Berlin}, 831-839 (1915).

{[}7{]} T. Damour and N. Deruelle, ``General relativistic celestial
mechanics of binary systems. I. The post-Newtonian motion,'' \emph{Ann.
Inst. Henri Poincare, Phys. Theor.} \textbf{43}, 107-132 (1985).

{[}8{]} T. Damour and G. Schaefer, ``Higher-order relativistic
periastron advances and binary pulsars,'' \emph{Nuovo Cimento B}
\textbf{101}, 127-176 (1988).

{[}9{]} L. E. Kidder, ``Coalescing binary systems of compact objects to
(post)\(^{5/2}\)-Newtonian order. V. Spin effects,'' \emph{Phys. Rev.~D}
\textbf{52}, 821-847 (1995).

{[}10{]} L. Blanchet, ``Gravitational radiation from post-Newtonian
sources and inspiralling compact binaries,'' \emph{Living Rev.~Relativ.}
\textbf{17}, 2 (2014).

\textbf{Gravitational radiation and waveforms}

{[}11{]} A. Einstein, ``Ueber Gravitationswellen,'' \emph{Sitzungsber.
Preuss. Akad. Wiss. Berlin}, 154-167 (1918).

{[}12{]} P. C. Peters and J. Mathews, ``Gravitational radiation from
point masses in a Keplerian orbit,'' \emph{Phys. Rev.} \textbf{131},
435-440 (1963).

{[}13{]} P. C. Peters, ``Gravitational radiation and the motion of two
point masses,'' \emph{Phys. Rev.} \textbf{136}, B1224-B1232 (1964).

{[}14{]} M. Maggiore, \emph{Gravitational Waves, Vol. 1: Theory and
Experiments} (Oxford University Press, 2008).

\textbf{Clocks and distances}

{[}15{]} L. D. Landau and E. M. Lifshitz, \emph{The Classical Theory of
Fields}, 4th ed., Course of Theoretical Physics Vol. 2 (Pergamon,
Oxford, 1975).

{[}16{]} T. Damour and N. Deruelle, ``General relativistic celestial
mechanics of binary systems. II. The post-Newtonian timing formula,''
\emph{Ann. Inst. Henri Poincare, Phys. Theor.} \textbf{44}, 263-292
(1986).

{[}17{]} I. I. Shapiro, ``Fourth test of general relativity,''
\emph{Phys. Rev.~Lett.} \textbf{13}, 789-791 (1964).

\textbf{Earlier examples of reading masses and distances from
relativistic orbital effects}

{[}18{]} J. H. Taylor and J. M. Weisberg, ``A new test of general
relativity: Gravitational radiation and the binary pulsar PSR 1913+16,''
\emph{Astrophys. J.} \textbf{253}, 908-920 (1982).

{[}19{]} T. Damour and J. H. Taylor, ``Strong-field tests of
relativistic gravity and binary pulsars,'' \emph{Phys. Rev.~D}
\textbf{45}, 1840-1868 (1992).

{[}20{]} C. Cutler and E. E. Flanagan, ``Gravitational waves from
merging compact binaries: How accurately can one extract the binary's
parameters from the inspiral waveform?,'' \emph{Phys. Rev.~D}
\textbf{49}, 2658-2697 (1994).

{[}21{]} B. F. Schutz, ``Determining the Hubble constant from
gravitational wave observations,'' \emph{Nature} \textbf{323}, 310-311
(1986).

\textbf{Numerical methods}

{[}22{]} K. Levenberg, ``A method for the solution of certain non-linear
problems in least squares,'' \emph{Q. Appl. Math.} \textbf{2}, 164-168
(1944).

{[}23{]} D. W. Marquardt, ``An algorithm for least-squares estimation of
nonlinear parameters,'' \emph{J. Soc. Ind. Appl. Math.} \textbf{11},
431-441 (1963).

{[}24{]} G. H. Golub and V. Pereyra, ``The differentiation of
pseudo-inverses and nonlinear least squares problems whose variables
separate,'' \emph{SIAM J. Numer. Anal.} \textbf{10}, 413-432 (1973).

\end{document}
